% GPT-6 Astra Pro; /root/transport; 22 November 2026 UTC.
% Actual terminal coefficient source and its quotient/corner maps.
\section{HD-F1-T08: the field of the literal four-dimensional tensor factor}
\label{sec:f1-terminal-tensor-field}

Retain the terminal natural source of HD-I23 and HD-R19:
\[
 G=S_s\times\Gamma,\qquad
 V=W\otimes_K M,\qquad W=\operatorname{Aug}(K^s),\quad
 s=r+1,\quad r\ge9,\quad\dim_K M=4.
 \tag{F1T8.1}
\]
The displayed product and coefficient isomorphism are the actual
ones constructed there. The representation is faithful and absolutely
irreducible, its exact character field is \(K\), and its specified
primitive Jacobian corner is a Weil sixfold of signature \((3,3)\)
and period rank nine in a still-missing determinant component.
Then
\[
                  K\in\{\mathbb Q(i),\mathbb Q(\sqrt{-2}),
                                     \mathbb Q(\sqrt{-7})\}.
 \tag{F1T8.2}
\]
The entire terminal tensor alternative is absent for the five fields
\[
 K=\mathbb Q(\sqrt{-d}),\qquad d\in\{11,19,43,67,163\}.
\]

Indeed the augmentation matrices are rational, and on the actual
subgroup \(\{1\}\times\Gamma\) the trace is
\[
          \operatorname{tr}_V(1,\gamma)=r\operatorname{tr}_M(\gamma).
 \tag{F1T8.3}
\]
Conversely every tensor trace is a rational augmentation trace
times \(\operatorname{tr}_M(\gamma)\). The two trace fields are
therefore equal to \(K\). Applying the already proved HD-R12 to
this literal finite subgroup of \(GL_4(K)\) gives
\[
             4\ge (d-1)/2.
 \tag{F1T8.4}
\]
For completeness, its exact field argument is as follows. A
nonrational trace supplies a nonrational irreducible factor \(f\)
of a finite-order matrix's characteristic polynomial over \(K\).
For a root \(\zeta_N\), restriction gives
\[
 \operatorname{Gal}(K(\zeta_N)/K)
   \simeq\operatorname{Gal}(\mathbb Q(\zeta_N)/
                              (K\cap\mathbb Q(\zeta_N))).
\]
If the intersection is \(\mathbb Q\), the factor is the rational
cyclotomic polynomial, a contradiction. Thus
\(K\subset\mathbb Q(\zeta_N)\) and \(\deg f=\varphi(N)/2\).
The ramified prime \(d\) in \(K\) ramifies in that actual field
tower, hence divides \(N\), giving \(\deg f\ge(d-1)/2\).
The corresponding summand is the image of the literal
Chinese-remainder projector \(u(\gamma)(\mu/f)(\gamma)\),
where \(u(\mu/f)+vf=1\) for the square-free minimal polynomial
\(\mu\). Its dimension is a positive multiple of \(\deg f\),
and cannot exceed four. The five lower bounds are
\(5,9,21,33,81\), proving the exclusion. This is an application
of HD-R12 to an actual coefficient factor, not a replacement
of the original Jacobian by a group with another character.

\section{HD-F1-T09: the terminal tensor rank is at least twenty}
\label{sec:f1-terminal-tensor-budget}

Every still-missing terminal source (F1T8.1) satisfies
\[
                r\ge20,\qquad m=4r\ge80,\qquad
                         1\le d_0\le4,
 \tag{F1T9.1}
\]
where \(d_0\) is the \(K\)-dimension of the parabolic cohomology
of the actual \(\Gamma\)-quotient coefficient system on the
original sphere. The proof uses no finite-group classification
beyond the actual product already established in HD-I23.

\paragraph{1. The actual smaller quotient and the zero case.}
Let \(C\) be the original connected \(G\)-cover of the sphere and
put
\[
 L_0=S_{s-1},\quad X=C/(L_0\times1),\quad Z=C/(S_s\times1),
 \quad Y=C/(L_0\times\Gamma).
 \tag{F1T9.2}
\]
There are actual quotient maps \(\alpha:C\to X\), of degree
\(|L_0|\), and \(f:X\to Z\), of degree \(s\);
\(X\to Y\) is a connected
\(\Gamma\)-cover. Diagonal inclusion and averaging on \(K^s\)
give the orthogonal coefficient decomposition
\[
 H^1_{\mathrm{par}}(Y,M)=
 H^1_{\mathrm{par}}(\mathbb P^1,M)\mathbin{\perp}
 H^1_{\mathrm{par}}(\mathbb P^1,W\otimes M).
 \tag{F1T9.3}
\]
The first pairing is scaled by \(s\); the dimensions are \(d_0\)
and six. At a sheet cycle of length \(\ell\), these actual maps
use \(\gamma^\ell\), so they identify the middle extensions at
the ramified points as well as the open local systems.

Carry the originally specified primitive coefficient vector by
the integer-cleared matrix unit of HD-I23 to \(w_0\otimes u\),
where \(w_0\) spans the \(L_0\)-fixed augmentation line. For the
primitive \(\Gamma\)-projector \(e_u\) on the quotient Jacobians,
the actual complementary operator is
\[
 \pi=e_u(1-s^{-1}f^*f_*),\qquad
 |L_0|^{-1}\alpha^*\pi\alpha_*
       =(\operatorname{av}_{L_0}-\operatorname{av}_{S_s})e_u.
 \tag{F1T9.4}
\]
These are the same actual coefficient and Jacobian maps as
HD-R16, (R16.15)--(R16.17); their proof does not use the
dimension three of its particular multiplicity space.
Integer-cleared images give isogenies of the prescribed
primitive factors. Quotient pullback scales the restricted
theta form by \(|L_0|\), and the change of primitive vector
has the positive rational multiplier
\(M_0^2h(v,v)/h(w_0\otimes u,w_0\otimes u)\), as in HD-I23.
On a rank-six Hermitian space each such scalar has sixth
power equal to the norm of its cube, so the ordinary
determinant class is retained. After finite parameter base
change these are flat Hodge maps and preserve the prescribed
period differential.

If \(d_0=0\), the connected image \(e_uJZ\) is zero. Since
the norm commutes with the \(\Gamma\)-projector,
\(f^*f_*e_u=0\), and (F1T9.4) is \(e_u\) on \(JX\).
Thus the entire rank-four quotient factor on \(Y\) is the
transported prescribed sixfold. Its Euler formula is
\[
             6=4(2g(Y)-2)+\sum e'_j .
 \tag{F1T9.5}
\]
Genus at least two is impossible; genus one has at most six
genuine marked coefficient branch values and period rank at
most six. Rank nine therefore forces \(Y=\mathbb P^1\).
The sum of its branch ranks is fourteen and there are at least
twelve such values, so one is a reflection. The already proved
HD-S36 for the Gaussian field and HD-G27 for the other fields
apply to this exact rank-four source, and give an already
supplied determinant. Consequently the still-missing source
has \(d_0\ge1\).

\paragraph{2. Exact local bounds with all phases retained.}
Write the original sphere tuple as \((\pi_j,\gamma_j)\), and put
\[
 c_j=\operatorname{rank}(P_{\pi_j}\otimes\gamma_j-I),\quad
 e_j=\operatorname{rank}(\gamma_j-I),\quad
 \delta_j=4-\max_\lambda\dim\ker(\gamma_j-\lambda I).
 \tag{F1T9.6}
\]
The original conditions and the quotient Euler formula give
\[
 \sum c_j=8r+6,\quad b\ge12,\quad c_j\ge4,\qquad
          \sum e_j=8+d_0\ge9.
 \tag{F1T9.7}
\]
The actual \(\gamma_j\)-tuple is irreducible and has product one.
The image-space proof of HD-R19, (R19.5), applied to this tuple
gives \(\sum\delta_j\ge5\). It permits arbitrary scalar phases
and requires no reflection-generation hypothesis.

For \(f_\ell(\gamma)=\dim\ker(\gamma^\ell-I)\), the actual
cycle-coordinate map followed by the diagonal/average splitting
gives
\[
 c=4r-\sum_{\text{cycles of }\pi}f_\ell(\gamma)+f_1(\gamma).
 \tag{F1T9.8}
\]
The following table exhausts the possible pairs \((e,\delta)\).
Its bounds are real lower bounds; \(c\) itself is an integer.
\[
\begin{array}{c|c|c}
 e&\delta&c\ \text{at least}\\\hline
0&0&4\\
4&0&2r-2\\
1&1&r\\
3&1&2r-1\\
4&1&(5r-3)/2\\
2&2&2r\\
3&2&(5r-1)/2\\
4&2&(8r-4)/3\\
3&3&3r\\
4&3&(16r-4)/5
\end{array}
 \tag{F1T9.9}
\]
Here the first row uses the original nonidentity coefficient
inertia. To verify all other rows, the bounds on \(f_\ell/\ell\)
are respectively \(2,3,2,3/2,2,3/2,4/3,1,4/5\).
For the two \(3/2\) rows, \(f_2\le3\), while \(\ell\ge3\)
gives \(f_\ell\le4\le3\ell/2\). For the \(4/3\) row there is
no eigenvalue one, the largest multiplicity is two, and
\(f_2\le2\); every \(\ell\ge3\) satisfies \(f_\ell\le4\).
When all four eigenvalues are distinct, \(f_\ell\le\ell\);
when one is also absent, \(f_\ell\le\min(4,\ell-1)\le4\ell/5\).
Substitute \(f_1=4-e\) and \(\sum\ell=r+1\) into (F1T9.8).

Two sharpenings will be useful. A nonidentity scalar has
\[
                  c\ge4(r-\lfloor(r+1)/2\rfloor).
 \tag{F1T9.10}
\]
For \((e,\delta)=(4,2)\), either the spectrum is
\((\omega,\omega,\omega^2,\omega^2)\), in which case
\[
                  c=4(r-\#\{\text{cycles with }3\mid\ell\})
                    \ge4(r-\lfloor(r+1)/3\rfloor),
 \tag{F1T9.11}
\]
or \(f_3\le3\), and then \(f_\ell\le\ell\) for every \(\ell\),
giving \(c\ge3r-1\). Both sharpenings follow on the actual
cycle eigenspaces. The scalar's inverse eigenvalue appears
at most once on a cycle and only at lengths divisible by its
order, which proves (F1T9.10).

\paragraph{3. Ranks nine through sixteen.}
For \(9\le r\le17\), every row of (F1T9.9) gives
\[
             c_j\ge4+\frac{r-5}{2}\delta_j+
                            \frac{r-3}{2}e_j.
 \tag{F1T9.12}
\]
For the \((4,2)\) row the difference is \((17-r)/3\);
for the \((4,3)\) row it is \((87-3r)/10\);
all other differences are \(0,2,3,4,5,8\) as read from
the table. Hence
\[
                    8r+6\ge 7r+22.
 \tag{F1T9.13}
\]
This excludes \(9\le r\le15\). At \(r=16\), equality
forces \(b=12,\sum\delta_j=5,\sum e_j=9\), with only
the rows \((0,0),(4,0),(1,1)\). Their numbers are six,
one and five respectively. The scalar row would have to
have \(c=30\), whereas (F1T9.10) gives \(c\ge32\).
Thus rank sixteen is excluded too.

\paragraph{4. Ranks seventeen through nineteen.}
The sharpenings make (F1T9.12) valid also for \(r=18,19\).
For \(r=17,18,19\), the total excess over (F1T9.13) is
respectively \(1,2,3\). Since increasing \(b\), \(\sum\delta_j\),
or \(\sum e_j\) would cost at least four, all three equal
\(12,5,9\).

At \(r=17\), the only rows with possible excess at most one
are \((0,0),(4,0),(1,1),(4,2)\). Write their nonidentity
\(\gamma\)-counts as \(z,a,p\) in the last three applicable
types, so
\[
                    a+2p=5,\qquad a+4z+4p=9.
 \tag{F1T9.14}
\]
Thus either \((z,p,a)=(1,0,5)\) or \((0,2,1)\).
None of the permitted rows can have excess one. The identity
and scalar costs are multiples of four. For a reflection
\(\gamma\), if \(a_\pi=\operatorname{rank}(P_\pi-I)\) and
\(t\) is the multiplicity of its nonunit inverse eigenvalue
on \(W\), then
\[
                         c=r+3a_\pi-t,\qquad t\le a_\pi,
 \tag{F1T9.15}
\]
so a nonzero excess is at least two. For the \((4,2)\) row,
the non-cubic alternative of (F1T9.11) has cost at least
fifty, while the cubic costs are multiples of four above
forty-four. This contradicts the total excess one.

At \(r=18\), the possible additional row \((3,1)\) has
excess at least two; call its count \(n\). The count equations
are \(n+2z+p=2\), and their minimum excess is
\(2n+2z+p\). Only \((n,z,p)=(0,1,0)\) and \((0,0,2)\)
can have total excess two, with respectively five and one
reflection entries. All local costs then attain their minima.
In the first case, the permutation tuple consists of six
transpositions, one product of nine disjoint transpositions
on nineteen letters, and identities. Its total permutation
index is fifteen. In the second, it consists of nine
transpositions and two permutations of cycle type \(3^6 1\);
its total index is thirty-three. But the actual connected
degree-nineteen quotient \(Y\to\mathbb P^1\) gives
\(\sum\operatorname{ind}(\pi_j)=2r+2g(Y)\ge36\).
Both cases are impossible.

At \(r=19\), the rows \((4,1)\) and \((4,3)\) may each
have excess three. If their counts are \(d,t\), subtraction
of the two count equations gives
\[
                         2n+4z+2p+3d+t=4.
\]
One such extra row forces the other, by parity and the
equation \(\sum\delta=5\); their combined excess is at least
six. Thus both are absent. The remaining equation is
\(n+2z+p=2\); excess at most three leaves only
\((n,z,p)=(0,1,0)\), five reflections, and six identity
\(\gamma\)-entries. The scalar cost is its minimum thirty-six,
and the identity costs are four, since their next costs
increase by four. The five reflection entries have total
excess three, so (F1T9.15) gives total permutation index at
most one among them. The scalar permutation has at most ten
cycles with eigenvalue \(-1\), and at its minimum is exactly
ten disjoint transpositions. Hence the full permutation
index is at most \(10+6+1=17\), whereas Riemann--Hurwitz
requires at least \(38\). This excludes rank nineteen and
proves \(r\ge20\).

\paragraph{5. A uniform upper bound on the discarded base summand.}
Let \(k\) count the nonidentity \(\gamma_j\), and let \(l\)
count those with \(e_j\ge2\). Formula (F1T9.9) gives \(c\ge r\)
for all nonidentity \(\gamma\), so
\[
             8r+6\ge k(r-4)+48,\qquad k\le7.
\]
For \(e\ge2\) it also gives \(c\ge2r-2\). Therefore
\[
             8r+6\ge r(k+l)+48-4k-2l.
\]
The alternative \(k+l\ge8\) would give at least \(8r+18\),
using \(k\le7\); hence \(k+l\le7\).
Now \(\sum e_j\le k+3l\le13\). Equality thirteen requires
three entries of rank four, one reflection, and all other
\(\gamma_j=1\). The three rank-four entries must have
projective codimensions whose sum is at least four.
The coordinate-minimal sorted triples are
\((0,1,3),(0,2,2),(1,1,2)\).
By (F1T9.9), after adding the reflection and the other
eight original inertias, their total costs are at least
\[
 \frac{87r+277}{10},\qquad
 \frac{25r+82}{3},\qquad
 \frac{26r+83}{3},
\]
respectively. Each exceeds \(8r+6\) for \(r\ge9\).
Thus \(\sum e_j\le12\), proving \(d_0\le4\).

\paragraph{6. Source and framework scope.}
The field input is the already proved HD-R12, with Milne's
original cyclotomic degree and ramification locator recorded
there; no new classification source is needed. The tuple
image-space inequality is HD-R19's complete elementary proof.
The fixed-space bounds, strengthened rank threshold and base
summand bound above are independent calculations on the actual
coefficient tuple. The accompanying rational-arithmetic script
checks the ten rows, their excesses, the endpoint count
equations and the numerical contradictions.

The maps (F1T9.2)--(F1T9.4) are finite quotient correspondences
and integer-cleared Jacobian homomorphisms. Write
\(n=f_*:JX\to JZ\) and \(j=f^*:JZ\to JX\).
Their normalized Poincare biextension identities give
\(\widehat n\lambda_Z=\lambda_Xj\) and
\(\lambda_{n^*\Theta_Z}=\widehat n\lambda_Zn\);
also \(j^*\Theta_X=s\Theta_Z\).
These are the restriction/trace polarizations above. They act on relative
Betti and de Rham realizations as flat paired maps after the
finite change fixing the group labels. Thus the same period
rank and full variational Hodge class of the prescribed
corner are retained. On \(K_0\), these proper and flat maps
use the actual pullback and derived pushforward of their
kernels; GRR realizes their cohomology action, and Adams
\(\psi^a\) multiplies codimension \(j\) by \(a^j\), as in F04.
The relative quotient curves and Jacobian graphs are finitely
presented algebraic families, giving their actual Hilbert and
Chow parameter maps. These exclusions supply no new detector
cycle or claim of semiregularity for an unspecified object.
They strictly reduce the natural-source alternatives in
T2-OTHER, while the no-moving-rank-at-most-four branch remains.
